\documentclass{cumcmthesis}
% 电子版提交时通常可使用：\documentclass[withoutpreface]{cumcmthesis}

% Template-First canonical layout lineage: introduced in v8.0.1 from the A196-inspired template work.
% This comment records provenance only; the current Skill release is declared by active release carriers.
% main.tex 只负责编排；固定一级骨架由 template_manifest.yaml 管理。
% 全局宏包与环境只在此处维护；正文 section 不重复加载宏包。
\usepackage{booktabs}
\usepackage{amsmath,amssymb,bm}
\usepackage{subcaption}
\usepackage{graphicx}
\usepackage{float}
\usepackage{algorithm}
\usepackage{algorithmic}
\usepackage{xcolor}
\usepackage{listings}
\usepackage{cleveref}
\usepackage[most]{tcolorbox}

% 保留用户参考模板中与官方 cumcmthesis 兼容的代码附录能力；不覆盖官方页面与标题格式。
\lstset{
  basicstyle=\ttfamily\small,
  keywordstyle=\color{blue!70!black}\bfseries,
  stringstyle=\color{red!65!black},
  commentstyle=\color{green!45!black},
  numbers=left,
  numberstyle=\tiny\color{gray},
  numbersep=6pt,
  frame=single,
  tabsize=4,
  captionpos=b,
  showstringspaces=false,
  breaklines=true
}

% cumcmthesis v2.6 自定义了旧版 \supercite；在加载 biblatex 前释放该命令。
\let\supercite\relax
\usepackage[backend=biber,style=numeric,sorting=none]{biblatex}
\addbibresource{references.bib}

% 命题按阿拉伯章节号编号。0--4 是默认正文阅读预算，不是数学硬上限。
% 短证明默认自然分段；只有明显多阶段证明才使用 enumerate。
\newtheorem{proposition}{命题}[section]
\renewcommand{\theproposition}{\arabic{section}.\arabic{proposition}}
\newenvironment{hskproposition}[1]
  {\begin{tcolorbox}[
      colback=white,
      colframe=black!72,
      boxrule=0.65pt,
      arc=1mm,
      outer arc=1mm,
      left=8pt,
      right=8pt,
      top=7pt,
      bottom=7pt,
      before skip=8pt,
      after skip=8pt,
      boxsep=0pt
    ]
    \begin{proposition}[#1]}
  {\end{proposition}\end{tcolorbox}}
\newenvironment{hskproof}
  {\par\smallskip\noindent\textbf{证明：}\par\nopagebreak\smallskip}
  {\hfill$\square$\par}

% 项目级可复用命令统一放在这里。
% 仅定义真正跨章节复用且语义稳定的命令；不要把一次性公式或局部变量包装成宏。

% 按当届题面和队伍信息修改，不在模板中固定年份、题号和小问数。
\newcommand{\HSKProblemCode}{A}
\newcommand{\HSKTeamNumber}{0000}
\newcommand{\HSKContestYear}{YYYY}
\newcommand{\HSKContestMonth}{9}
\newcommand{\HSKContestDay}{1}

\title{XXX问题的XXX建模与分析}
\tihao{\HSKProblemCode}
\baominghao{\HSKTeamNumber}
\schoolname{学校名称}
\membera{成员A}
\memberb{成员B}
\memberc{成员C}
\supervisor{}
\yearinput{\HSKContestYear}
\monthinput{\HSKContestMonth}
\dayinput{\HSKContestDay}


\begin{document}
\maketitle

\begin{abstract}
% 第一段通常 2--3 句：研究对象 + 核心矛盾 + 全文统一路线。不要在第一段堆全部小问方法清单。
……

% 随后按问覆盖“任务/对象 + 模型 + 目标/关键条件或约束 + 算法/求解方法
% + 决定性结果 + 真实检验证据（若有）+ 直接结论”。这是内容链，不是固定句式；
% 不得为了补齐链条虚构敏感性、鲁棒性或替代算法检验，也避免每段都复制“针对问题X，本文……”。
% 核心答案精度与 Numeric Profile、正文关键结果表保持一致。
……

\keywords{关键词1\quad 关键词2\quad 关键词3\quad 关键词4}
\end{abstract}


\section{问题重述}

\subsection{问题背景}
% 默认一个自然段。利用题目给出的现实背景和领域矛盾，自然收束到待解决的问题。
% 不另起一段介绍“全文结构”“后文安排”，也不机械扩写宏大背景。

\subsection{问题提出}
% 先抽取对象、关键条件/范围/量词/单位与待求输出，再按逻辑依赖重新组织。
% 不照抄原题，不沿原题句序逐句换同义词；同时不得为了改写漏掉决定题意的条件。
% 不提前写模型名、公式、算法、结果，也不把模型假设改写成题目已知条件。
\noindent\textbf{问题一：}……

\noindent\textbf{问题二：}……

% 按题目实际小问继续，不固定小问数。

\section{问题分析}

% 参考优秀论文 A196：问题分析按小问逐项展开。每问通常使用逻辑连续的自然段，
% 形成“当前对象/条件 → 真正难点 → 数学抓手 → 建模转化 → 准备建立的结构
% → 与前问关系”的承接，不写成散乱短句、名词清单或软件流水线；
% 不提前展开完整公式、算法细节和最终结果。

\subsection{问题一的分析}
% 说明问题一的研究对象、关键困难、所需判定/估计/优化量，以及为什么需要后续模型准备或本问模型。

\subsection{问题二的分析}
% 若依赖问题一，明确继承了什么结果、还新增什么变量/目标/约束；不要重新复述问题一全部机制。

% 按实际小问数量复制，不固定五问：
% \subsection{问题三的分析}
% \subsection{问题四的分析}
% \subsection{问题五的分析}

% 本章不写“先读取数据、再建模、再求解、再绘图”的软件流水线；
% 应从题目对象和数学困难出发形成各问之间的递进关系。

\section{模型假设}
% 只保留会改变变量、约束、目标、状态、分布、近似误差或适用边界的假设；
% 每项应能说明来源/必要性、影响的模型部分，以及必要时的失效影响或检验位置。
% 题面事实、单位约定和程序正确性不包装成假设。
\begin{enumerate}
    \item ……。
    \item ……。
    \item ……。
\end{enumerate}

\section{符号说明}
% 只列全文高频符号。局部变量在首次出现处定义；场景/模型/方案优先使用短上标，避免长文本和多层下标。
\begin{table}[H]
    \centering
    \caption{主要符号说明}
    \label{tab:symbols}
    \begin{tabular}{cccc}
        \toprule
        符号 & 含义 & 单位 & 类型 \\
        \midrule
        $x_i$ & 元素/位置索引示例 & -- & 变量 \\
        $I^{(s)}$ & 第 $s$ 个模型对应结果 & -- & 状态量 \\
        $\theta$ & 参数示例 & -- & 参数 \\
        \bottomrule
    \end{tabular}
\end{table}


% 数据说明/项目级预处理是条件章节。仅当确有共享数据方法需要时取消下一行注释。
% % preprocessing_decision=not_needed 时删除 main.tex 中对本文件的 \input；
% question_local 时把必要变换放入对应问题；只有 project_level 才保留实质公共数据预处理章节。
\section{数据预处理}

\subsection{数据问题与处理必要性}
% 说明真实数据问题；不得为了形式完整编造清洗、插值、标准化。

\subsection{处理方法与验证}
% project_level 时给公式、参数依据、验证、处理前后证据和后续模型接口。


% 参考优秀论文的章节逻辑：若两问及以上共享实质机理、轨迹、几何关系、公共判定式或状态方程，
% 在正式进入问题一前统一写“模型准备”；仅满足该条件时取消下一行注释。
% \section{模型准备}

% 本章参考优秀论文 A196 的章节组织，用于统一承接“两问及以上共享”的基础机理。
% 若当前赛题不存在跨问共享的实质模型结构，应删除本章，而不是为了编号完整强制保留。
% 本章不写各问最终优化目标、solver、主结果和验证结论。

\subsection{共享机理与基础关系}
% 只定义后续多问都会复用的对象、坐标系、运动/状态关系、守恒关系、公共参数或几何结构。
% 每个核心关系都应说明来源、现实含义以及它将被哪些问题调用。

% 示例占位：
% 设公共状态向量为 $\mathbf{s}(t)$，则基础演化关系可写为
% \begin{equation}
%     \dot{\mathbf{s}}(t)=F\bigl(\mathbf{s}(t),\boldsymbol{\theta}\bigr).
%     \label{eq:shared-state}
% \end{equation}
% 式~\eqref{eq:shared-state} 只表示共享基础关系；正式项目必须替换为当前题目的真实机理。

\subsection{公共判定关系或约束来源}
% 统一定义后续各问反复使用的事件判定、可行域、几何关系、阈值、边界或评价指标。
% 如果判定关系只服务一个问题，则不要放在“模型准备”，应写回对应问题章节。

% 复杂机理题可在此配置 1--2 张核心机理图，帮助解释变量、边界和判定条件；
% 数据处理本身不属于本章，项目级预处理应使用条件章节 05_data.tex。


% 每个问题独立一级章，统一采用“问题X模型建立及求解”。
\section{问题一模型建立及求解}

% v8.0.1 A196-inspired question chapter.
% 复杂问题默认形成“模型建立 → 模型求解 → 求解结果 → 结果的分析与验证”的局部闭环；
% 这些是功能阶段，不是不可修改的标题。若更专业的题目专属标题能更准确暴露数学对象，可直接替换。
% 简单解析题或直接计算题可以合并、删减非必要阶段。

\subsection{模型建立}
% 建议顺序：对象与当前缺口 → 变量/状态 → 机理或定义 → 目标/判定关系 → 约束 → 可计算模型。
% 优化类问题可按需设置三级标题：决策变量、目标函数、约束条件、模型汇总；
% 机理类问题则应按真实对象设置专业标题，不机械套“决策变量”。

设……，由前述共享关系/本问条件可得
\begin{equation}
    F(x,\theta)=0,
    \label{eq:q1-core-relation}
\end{equation}
式~\eqref{eq:q1-core-relation} 给出……，并作为后续求解的直接输入。

% 命题只在承担等价变换、单调性、可行性、降维或边界判定等真实作用时使用。
% \begin{hskproposition}{……}
% \label{prop:q1-example}
% ……
% \begin{hskproof}
% ……
% \end{hskproof}
% \end{hskproposition}

% 对复杂优化模型，模型建立末尾可进行核心模型收束。
% 目标函数必须单独展示，不进入约束大括号：
\begin{equation}
    \min_{\mathbf{x}} f(\mathbf{x})
    \label{eq:q1-objective}
\end{equation}

\begin{equation}
\text{s.t.}\quad
\left\{
\begin{aligned}
    g_i(\mathbf{x}) &\le 0, \quad i=1,\ldots,m, \\
    h_j(\mathbf{x}) &= 0, \quad j=1,\ldots,n, \\
    \mathbf{x} &\in \Omega .
\end{aligned}
\right.
\label{eq:q1-constraints}
\end{equation}
% 非优化题必须删除上述 optimization 示例，改成真实状态方程、概率映射、边界条件或解析关系。
% “模型汇总”是模型建立末尾的收束动作，不默认单列“核心模型汇总”二级标题。

\subsection{模型求解}
% 从“最终模型已经化成什么计算问题”开始，而不是先介绍算法百科。
% 写清：剩余计算困难 → solver/离散/搜索/分解方法为什么适配 → 关键参数、初值、精度、终止条件 → 输出映射。
% 若循环、分支、修复、接受/拒绝或终止逻辑本身是关键证据，可在解释输入与原因后再给算法步骤/伪代码。

\subsection{求解结果}
% 参考优秀论文的局部证据链：主结果立即在本问展示，图/表之后紧跟决定性数值、趋势和机制解释。
% 不要连续堆多张图后再统一解释；本节末尾应自然、直接回答问题一。

\subsection{结果的分析与验证}
% 当存在真实风险需要检查时启用：离散精度、参数敏感性、鲁棒性、多初值/多 seed、边界扰动、
% 替代算法/模型、残差/误差区间、外样本等。每项验证说明：检验什么风险 → 改变什么 → 指标 → 结果变化 → 支持边界。
% 若问题一属于简单解析/直接计算且没有实质验证必要，可删除整个小节。

% 按实际小问复制为 07_question2.tex、08_question3.tex 等。
% 后问继承前文已经定义的共享模型，只写新增变量、目标、约束、算法和证据，不从头复制模型准备。

% Q2/Q3 示例文件已经提供，按实际小问数量启用；Q4/Q5 可复制 Q3 后改名：
% \section{问题二模型建立及求解}

% 这是复杂优化类问题的第二个示例章节，可按当前题目直接复制修改。
% 若问题二只是问题一的解析延伸，应删减不必要的小节，不为了模板对称强行保留完整四段。

\subsection{优化模型建立}

% 可按需设置三级标题，例如：
% \subsubsection{决策变量定义}
% \subsubsection{目标函数设置}
% \subsubsection{约束条件分析}
% \subsubsection{模型收束}
% 三级标题必须服务真实数学对象，不要求固定名称或固定数量。

设决策向量为 $\mathbf{x}$，目标函数为 $f(\mathbf{x})$。则本问优化目标写为
\begin{equation}
    \max_{\mathbf{x}} f(\mathbf{x}).
    \label{eq:q2-objective}
\end{equation}

约束条件统一写为
\begin{equation}
\text{s.t.}\quad
\left\{
\begin{aligned}
    g_i(\mathbf{x}) &\le 0, \quad i=1,\ldots,m,\\
    h_j(\mathbf{x}) &= 0, \quad j=1,\ldots,n,\\
    \mathbf{x} &\in \Omega.
\end{aligned}
\right.
\label{eq:q2-constraints}
\end{equation}

% 正式项目必须用真实公式替换示例；目标函数保持在约束大括号外。

\subsection{优化模型求解}
% 说明模型结构、搜索空间/非凸性/离散性/耦合约束等实际困难，随后解释 solver 或启发式为何适配。
% 多阶段算法应说明每一步输入、输出及其如何进入下一步；不写通用算法百科。

\subsection{求解结果}
% 给出最优/推荐方案、核心变量、目标值和必要图表，并紧邻解释其对问题二的含义。

\subsection{结果的分析与验证}
% 根据本问真实风险设置敏感性、鲁棒性、替代算法、边界扰动、局部邻域、最优性间隙或多启动一致性等验证。
% 验证的目标是界定当前结论的可信范围，而不是形式化补一个“敏感性分析”小节。

% \section{问题三模型建立及求解}

% 后续问题模板：先说明真实继承，再说明问题三新增的对象、条件或数学困难。
% 不复制问题一/二的完整推导；若问题三独立，直接说明新的研究对象和数学缺口。

\subsection{模型建立与扩展}
% 建议顺序：继承的 current 关系 → 新增变量/状态 → 新目标或判据 → 新约束来源 → 最终可计算模型。
% 决定性推导、边界和约束不得因“前文已有模型”而省略；未变化部分用准确引用承接。
% 优化类仍应先说明决策对象和目标函数现实含义，并把目标函数放在约束大括号外。

\subsection{模型求解}
% 从问题三最终模型的计算结构或新增困难进入 solver。
% 写清继承/更换求解方法的原因、本题化编码、约束处理、初值/参数/精度/终止条件和输出映射。
% 算法步骤或伪代码只有在真实多阶段、循环、分支或筛选逻辑需要展示时才启用。

\subsection{求解结果}
% 给出问题三的 current 高精度核心答案、推荐方案或判定，并在核心图表邻近解释决定性特征。
% 结果段应直接回答问题三；不能只停在算法收敛、目标值或“结果如图所示”。

\subsection{结果的分析与验证}
% 仅在存在真实风险和 current 证据时启用。验证前先说明主结果已回答什么、仍受什么因素影响。
% 每项检验闭合：风险/主张 → 改变什么 → 保持什么 → 指标 → 变化范围 → 结论及有效边界。
% 若问题三简单且没有独立验证必要，删除整个小节。

% \input{sections/09_question4}
% \input{sections/10_question5}

\section{模型的评价与推广}

% 参考优秀论文 A196 的末章组织：评价必须对应前文已经得到的模型、结果和验证证据，
% 不是统一套用“优点很多、推广性强”等空泛表述。

\subsection{模型的优点}
% 数量由实际模型决定；每一点都说明“哪一项结构/机制/数值证据”带来了什么具体优势。
\begin{enumerate}
    \item ……。
    \item ……。
\end{enumerate}

\subsection{模型的不足与改进方向}
% 不足要说明会影响哪个结果、参数范围或适用边界；改进方向必须与该不足一一对应。
\begin{enumerate}
    \item ……。
    \item ……。
\end{enumerate}

\subsection{模型的推广}
% 选写。只有存在明确可迁移对象时保留，说明哪些模型结构可直接继承、哪些参数/边界需要重新标定。

% 中文国赛默认不单设“结论”一级章；各问题的最终回答应已经在对应“求解结果/结果分析”中闭环。


% AI 工具使用声明是条件合规槽位。只有当届规则已核验、当前项目真实 AI 使用事实已确认，
% 且披露在当前项目中适用时，先据实填写 sections/10_ai_tool_statement.tex，再取消下一行注释。
% % AI 工具使用声明（条件合规槽位）
% 本文件只提供结构位置与填写边界，不陈述任何具体参赛队的 AI 使用事实。
%
% 默认状态：hsk_main.tex 中对应 \input 保持注释，不进入最终 PDF。
% 仅当以下条件均已满足时才启用并填写：
% 1. 当前竞赛当届规则已经核验，规则来源为
%    config/competition_profiles.yaml#profiles.<name>.edition_rules.ai_disclosure；
% 2. 当前项目的真实 AI 使用事实已经由用户/参赛队确认；
% 3. 当前规则与真实事实共同表明需要或适合在论文中披露。
%
% 启用后必须把本文件改写为当前项目的真实声明；不得把模板说明原样提交，
% 不得默认声称使用了 AI，也不得根据论文文风推断 AI 使用情况。
% 最终一致性由 modules/06_review_delivery.md#Final-Submission-Compliance-Evidence-Sweep
% 中的 ai_disclosure 检查族核对。
%
% 填写结构示意（以下均为注释，不会出现在 PDF）：
% \section*{AI工具使用声明}
% <仅填写已核验规则要求/允许披露且参赛队确认属实的工具、用途、范围与责任说明。>


\printbibliography[title={参考文献}]

\begin{appendices}

\section{核心代码与复现说明}
% 正文只保留必要算法流程；完整 Python/MATLAB 可放附录或独立附件，并说明输入、功能和输出。

\section{补充推导与图表}
% 技术证明超过正文安全长度时放在此处；仅放正文不宜展开但仍有证据价值的图表。

\end{appendices}


\end{document}
